\documentclass{beamer}
\usepackage[utf8]{inputenc}
\usepackage[spanish]{babel}
\usepackage{hyperref}
\usepackage{graphicx}
\usepackage{wrapfig}
\usepackage[square,numbers]{natbib}
%\bibliographystyle{unsrtnat}

\hypersetup{
    colorlinks=true,
    linkcolor=blue,
    filecolor=magenta,      
    urlcolor=cyan,
}

\usetheme{Madrid}
\usecolortheme{default}

\title[ASIGNATURA] %optional
{Nombre de la asignatura}
\author[Autor] % (optional)
{Nombre del autor\inst{1}}
\institute[UniMinuto] % (optional)
{
  \inst{1}%
  Facultad de Ingeniería\\
  Universidad Minuto de Dios
%  \and
%  \inst{2}%
%  Faculty of Chemistry\\
%  Very Famous University
}
\date[CPCXXXX-YYYY-Nombre @Año] % (optional)
{AAAA-semestre\\Horario\\Salón de clase}
\logo{\includegraphics[height=1.5cm]{Logo_UniMinuto.png}}

\AtBeginSection[]
{
  \begin{frame}
    \frametitle{Contenido}
    \tableofcontents[currentsection]
  \end{frame}
}

\begin{document}
\frame{\titlepage}
\section{Nombre de la sección}
\begin{frame}
\frametitle{Nombre de la diapositiva}
Contenido de la diapositiva. Puede tener viñetas sin numerar
\begin{itemize}
    \item \emph{Se puede hacer un llamado de atención},
    \item O normalito,
     \begin{itemize}
         \item Se pueden colocar viñetas de segundo y n-nivel
     \end{itemize}
    \item Programación dirigida por eventos,
    \begin{enumerate}
        \item Numeradas también
    \end{enumerate}
\end{itemize}
\begin{enumerate}
    \item O al revés
    \begin{itemize}
        \item La numeración es en el primer nivel
    \end{itemize}
\end{enumerate}
\end{frame}

\section{Otra sección - Esto aparece en el gestor de contenidos}
\begin{frame}
\frametitle{El nombre de la diapositiva}
Se permiten bloques de llamado de atención
\begin{alertblock}{No todo es blanco o negro}
A veces se puede llamar la atención.
\end{alertblock}
\begin{columns}
\column{0.5\textwidth}
Y qué tal si se requiere texto en columnas?
\begin{enumerate}
    \item Desarrollo más rápido,
\end{enumerate}
    
\column{0.5\textwidth}
Segunda columna\\
Comenzamos con las imágenes:
\begin{figure}
    \centering
    \includegraphics[scale=0.2]{Logo_UniMinuto.png}
    \caption{El nombre de la figura}
    \label{fig:paraHacerLlamadosEnElCuerpo}
\end{figure}
\end{columns}
\end{frame}

\begin{frame}{Dentro de la misma sección, otro título}
Si se requieren vínculos:
\href{https://bit.ly/31Xqmaf}{Lenguajes}
\end{frame}

\subsection{Una subsección}
\begin{frame}{Siempre se nomenclan las diapositivas}
\href{https://www.youtube.com/watch?v=-SADbPS8UgA}{Arquitectura von Neumann}
\end{frame}

\section{Bibliografía}
\begin{frame}{También hay espacio para las referencias}
\bibliographystyle{abbrvnat}
\begin{thebibliography}{9}
\bibitem{Fernandez-Pena2011} 
Fernández-Peña, Juan M. 
\textit{\href{https://www.uv.mx/personal/jfernandez/files/2010/07/8_Calidad.pdf}{Modelos de calidad de software}}. 2011

\bibitem{Pressman2010} 
Pressman, Roger S. 
\textit{\href{http://cotana.informatica.edu.bo/downloads/ld-Ingenieria.de.software.enfoque.practico.7ed.Pressman.PDF}{Ingeniería de Software: Un enfoque práctico 7a ed}}. Mc Graw-Hill, 2010

\bibitem{Callejas2017} 
Callejas-Cuervo, Mauro, Alarcón-Aldana, Andrea C. y Álvarez-Carreño, Ana M.
\textit{\href{hhttp://www.scielo.org.co/scielo.php?script=sci_abstract&pid=S1900-38032017000100236&lng=en&nrm=iso&tlng=es}{Modelos de calidad del software, un estado del arte}}. 2017

\bibitem{Fuertes2008}
Fuertes-Castro, José L.
\textit{\href{http://www.iimv.org/iimv-wp-1-0/resources/uploads/2015/01/JoseLuisFuertesCalidaddelSoftware.pdf}{Calidad de Software}}. 2008
\end{thebibliography}
\end{frame}
\begin{frame}
\frametitle{Gracias}
\begin{figure}[h]
%\caption{Gracias}
\centering
\includegraphics[scale=0.6]{agradecimiento.jpg}
\end{figure}
\end{frame}

\end{document}